\documentclass[9pt]{beamer}
\usepackage[latin1]{inputenc}
\usepackage{multicol,multirow,dsfont}
\usepackage{graphicx}
\usepackage{multimedia}
\usepackage{color}
\usepackage{fancybox}
\usepackage{algorithm}
\usepackage{algorithmic}
%\usepackage{pifont}
\usetheme{Darmstadt}
\setbeamertemplate{footline} %customizing footline
{
\hbox{%
\begin{beamercolorbox}[wd=.47\paperwidth,ht=2.25ex,dp=1ex,leftskip=.3cm,rightskip=.3cm plus1fil]{author in head/foot}%
\insertframenumber{}/\inserttotalframenumber\hfill
\usebeamerfont{author in head/foot}\insertshortauthor
\end{beamercolorbox}%
\begin{beamercolorbox}[wd=.54\paperwidth,ht=2.25ex,dp=1ex,leftskip=.3cm,rightskip=.3cm ]{title in head/foot}%
\usebeamerfont{title in head/foot}\insertshorttitle
\end{beamercolorbox}%
}
\vskip0pt}

\newenvironment<>{varblock}[2][.9\textwidth]{%
\setlength{\textwidth}{#1}
\begin{actionenv}#3%
\def\insertblocktitle{#2}%
\par%
\usebeamertemplate{block begin}}
{\par%
\usebeamertemplate{block end}%
\end{actionenv}}

\newtheorem{DEF}{Definition}
\newtheorem{HYP}{Hypothesis}
\newtheorem{THEO}{Theorem}
\newtheorem{PROP}{Proposition}
\newtheorem{LEMM}{Lemma}
\newtheorem{CONJ}{Conjecture}
\newtheorem{REM}{Remark}
\newtheorem{EX}{Example}
\newtheorem{PRO}{Propriety}

\def\R{{\ensuremath{\mathrm{I\!R}}}}
\def\N{{\ensuremath{\mathrm{I\!N}}}}
\def\ab{\textbf{\emph{a}}}
\def\xb{\textbf{x}}
\def\sb{\textbf{s}}
\def\vb{\textbf{v}}

\definecolor{vert}{rgb}{0,0.7,0}
\definecolor{vert1}{rgb}{0,0.6,0}
\definecolor{rouge2}{rgb}{0.4,0,0}
\definecolor{bleuclair}{rgb}{0.9,0.9,1}
\definecolor{lightgreen}{rgb}{.5,1,.5}


%\setbeamertemplate{navigation symbols}{}

\title[]{Imperfect Maintenance and Gamma Process}

\author[]{Franck Corset*, Mitra Fouladirad**, Christian Paroissin*** \medskip}

\institute[UTT/LPPA]{
   \footnotesize{* Universit� Grenoble Alpes LJK\\
   **Universit\'{e} de Technologie de Troyes, LIST3N \\
***Universit\'e de Pau et des Pays de l'Adour, LMA }}

\date[UGA/UTT/LPPA, France]{\small MIMAR 2021}

\begin{document}
\frame{
\titlepage}

%\AtBeginSection[] {
%\begin{frame}<beamer>
%\frametitle{Summary}
%\tableofcontents[currentsection,hideothersubsections]
%\end{frame}
%}

%oooooooooooooooooooooooooooooooooooooooooooooooooooooooooooooooooooooo



%oooooooooooooooooooooooooooooooooooooooooooooooooooooooooooooooooooooo
%\section{}
%\frame{
%\frametitle{Plan de l'expos\'{e}}
%\tableofcontents
%}

%\section{Problem statement}
%
%\frame{
%\frametitle{}
%\begin{itemize}
%\item 
%\end{itemize}
%}
%%%%%%%%%%%%%%%%%%%%%%%%%%%%%%%%%%%%%%%%%
\frame{
\frametitle{Problem statement}
\begin{itemize}
\item Deteriorating system
\item Periodic inspections
\item Condition-based maintenance
\item AGAN corrective maintenance
\item Imperfect preventive maintenance
\end{itemize}
}
%%%%%%%%%%%%%%%%%%%%%%%%%%%%%%%%%%%%%%%%%

\frame{
\frametitle{Deterioration model}
The deterioration process is a Gamma process, $\{X_t,t\geq 0\}$, with shape function $a(t)=\alpha t^\beta$ with  $\alpha>0$ and $\beta>0$  and scale parameter $b>0$.

The Gamma process $\{X_t,t\geq 0\}$ is a continuous-time stochastic process such that:
\begin{itemize}
	\item $X_0 = 0$ with probability $1$
	\item $X_t-X_s\sim G(a(t)-a(s),b)$ for all $t>s\geq0$, with density function $\forall x\in \mathbb R^+$:
	$$
f(x) = \dfrac{b^{a(t)-a(s)}}{\Gamma(a(t)-a(s))}x^{a(t)-a(s)-1}e^{-bx}
	$$
	\item $\{X_t,t\geq 0\}$ has independent non overlapping increments 
\end{itemize}
}
%%%%%%%%%%%%%%%%%%%%%%%%%%%%%%%%%%%%%%%%%
\frame{
\frametitle{Maintenance policy}
\begin{itemize}
	\item $\{Y_t,t\geq 0\}$  the stochastic process of the maintained system. 
	\item Inspection times $(t_j)_{j\in \N}$, 
	\item Periodic inspections  $t_j=j\tau$, $j\in \N$. 
\end{itemize}
	Maintenance decision rule
\begin{itemize}
	\item if $Y_{t_j^-} \geq M$ then a AGAN CM $Y_{{t_j}^+} =0$,
	\item if $L\leq Y_{t_j^-} < M$ imperfect PM,  $Y_{{t_j}^+} = (1-\rho)Y_{{t_j}^-}$,
	\item if $Y_{t_j^-} < L$ no action is performed,
\end{itemize}
where $L$ and $M$ are fixed thresholds ($0<L<M$).
}
%%%%%%%%%%%%%%%%%%%%%%%%%%%%%%%%%%%%%%%%%
%%%%%%%%%%%%%%%%%%%%%%%%%%%%%%%%%%%%%%%%%
%\frame{
%\frametitle{}
%After the first inspection without replacement, we have $\forall t\in [\tau,2\tau)$, 
%$$
%Y_t = Y_\tau + X_t^{(2)} - X_\tau^{(2)} = (1-\rho)^{U_1} X_\tau^{(1)}+X_t^{(2)} - X_\tau^{(2)} 
%$$
%At the second inspection at time $2\tau$, we have
%\begin{itemize}
%	\item if $Y_{2\tau^-} <L$, no action is performed, and $Y_{2\tau} = Y_{2\tau^-} = (1-\rho)^{U_1} X_\tau^{(1)}+X_{2\tau}^{(2)} - X_\tau^{(2)}$.
%	\item if $L \leq Y_{2\tau^-} <M$, an imperfect PM ($ARD_1$) is performed and then the degradation level is reduced of $\rho (X_{2\tau}^{(2)}-X_\tau^{(2)})$ and thus 
%	$$
%	Y_{2\tau} = Y_{2\tau^{-}} -\rho (X_{2\tau^-}^{(2)}-X_\tau^{(2)}) = (1-\rho)^{U_1} X_\tau^{(1)}+(1-\rho)(X_{2\tau^-}^{(2)}-X_\tau^{(2)})
%	$$
%	\item if $Y_{2\tau^-} >M$ then the system is replaced.
%\end{itemize}
%}
\frame{
\frametitle{}
By induction, we can conclude that after $i$ inspections without replacement, we have $\forall t\in [i\tau,(i+1)\tau)$
$$
Y_t = Y_{i\tau} + (X_t^{(i+1)}-X_{i\tau}^{(i+1)})
$$
and if we denote $i=\lfloor t/\tau\rfloor$, then we get
$$
Y_t = \displaystyle\sum_{j=1}^{i} (1-\rho)^{1-U_j}(X_{j\tau}^{(j)}-X_{(j-1)\tau}^{(j)}) + X_t^{(i+1)}-X_{i\tau}^{(i+1)}
$$
Finally, we obtain
\begin{align}
Y_{(i+1)\tau} & = Y_{i\tau} + (1-\rho)^{U_{i+1}}(X_{(i+1)\tau}^{(i+1)}-X_{i\tau}^{(i+1)})\nonumber\\
& = \displaystyle\sum_{j=1}^{i+1} (1-\rho)^{U_{j}} (X_{j\tau}^{(j)}-X_{(j-1)\tau}^{(j)})
\end{align}

}
%%%%%%%%%%%%%%%%%%%%%%%%%%%%%%%%%%%%%%%%%
\frame{
\frametitle{Aim}
\begin{itemize}
	\item  Maintenance policy evaluation : cost criteria
	\item  Maintenance parameter optimisation
	\item  Practical situation: unknown deterioration parameters
	\end{itemize}
}
%%%%%%%%%%%%%%%%%%%%%%%%%%%%%%%%%%%%%%%%%
\frame{
\frametitle{Maintenance Cost evaluation}
Let $C_I$, $C_P$ and $C_C$ denote respectively the inspection cost, PM cost and CM cost.
The maintenance cost at time $t$ is given as follows:
$$C(t)=C_I [\frac{t}{\tau}] +C_P \sum_{i=1}^{[t/\tau]}{\mathds 1}_{L\leq Y_{i\tau^-}< M}+ C_C \sum_{i=1}^{[t/\tau]}{\mathds 1}_{ Y_{i\tau^-}\geq M}$$
As the corrective maintenances are perfect (replacement), by the renewal theory we can compute the long run average maintenance cost as follows:
\begin{align*}\label{eq-Ct-ECT}
\lim_{t\to\infty}\frac{C(t)}{t}&=\frac{{\mathds E}(C(T))}{{\mathds E}(T)}= %\frac{ {\mathds E}(C_C + [T/\tau] C_I + \sum_{i=1}^ {T/\tau} {\mathds 1}_{L \leq Y_{i \tau^-}< M}  C_p)
%}{{\mathds E}(T)}\nonumber\\
&\frac{C_c+ {\mathds E}([T/\tau]) C_i+ \sum_{i=1}^ {T/\tau} {\mathds P}({L \leq Y_{i \tau}< M}) C_P}{{\mathds E}(T)}
\end{align*}
where $T$ is the length of a renewal cycle which means the period between two corrective maintenance:
$$
T = \inf \{t>0,\,Y_t\geq M\}
$$

}
%%%%%%%%%%%%%%%%%%%%%%%%%%%%%%%%%%%%%%%%%
\frame{
\frametitle{}
To bypass this problem, the markovian renewal property of the maintained degradation process between two inspection is used to calculate the long run average maintenance cost as follows:
\begin{align*}
\lim_{t\to\infty}\frac{C(t)}{t}&=\frac{{\mathds E}_{\pi}(C(\tau))}{{\mathds E}_{\pi}(\tau)}
%\frac{ C_I +  C_p{\mathds P}_{\pi}(L \leq Y_{ \tau^-}< M) +C_C {\mathds P}_{\pi}( Y_{ \tau^-}\geq M)
%}{{\mathds E}_{\pi}(\tau)}\nonumber\\
 = \frac{ C_I +  C_p{\mathds P}_{\pi}(L \leq Y_{ \tau^-}< M) +C_C {\mathds P}_{\pi}( Y_{ \tau^-}\geq M)}{\tau}
\end{align*}
where ${\mathds E}_{\pi}$ and ${\mathds P}_{\pi}$  are the expectation and the probability under the steady-state distribution $\pi$ respectively.
\begin{equation}\label{eq-Pi_1}
{\mathds P}_{\pi}(L \leq Y_{ \tau^-}< M)=\int_{0}^M \int_L^M f(y-x)\pi(dx)\,dy 
\end{equation}
\begin{equation}\label{eq-Pi_2}
{\mathds P}_{\pi}( Y_{ \tau^-}\geq M)=\int_{0}^{+\infty} \int_M^{+\infty} f(y-x)\pi(dx)\,dy 
\end{equation}
}

%%%%%%%%%%%%%%%%%%%%%%%%%%%%%%%%%%%%%%%%%


\frame{
\frametitle{Stationary law of the Maintained System}
Let consider the Markov renewal cycle $[t_{i-1}^-,t_i^-]$, $x$ and $y$ respectively the deterioration level of maintained system at the beginning and the end of the cycle, the steady-state distribution $\pi$ is the solution of :
\begin{equation}
    \pi(y) = \int_0^{+\infty} \pi(x) p(y|x)\,dx
\end{equation}
where $p(y|x)$ is the transition law from the state $x$ to the state $y$. This transition law is equal to
\begin{equation}
    p(y|x) = f(y-x){\mathds 1}_{x <L}  +f(y-(1-\rho)x){\mathds 1}_{L\leq x <M}  +f(y){\mathds 1}_{x >M} 
\end{equation}
where $f$ is the density function of gamma distribution :
$$
f(x) = \dfrac{\beta^{\alpha (t_i^\beta - t_{i-1}^\beta )}}{\Gamma(\alpha (t_i^\beta - t_{i-1}^\beta ))}x^{\alpha (t_i^\beta - t_{i-1}^\beta )-1}e^{-b x} {\mathds 1}_{\R^+}(x)
$$

}
%%%%%%%%%%%%%%
\frame{
\frametitle{Steady-state distribution of the Maintained System}
Thus, we have
\begin{align*}
\label{TransitionLaw}
      \pi(y) = \int_0^{\min(L,y)} \pi(x) f(y-x)\,dx &+ \int_{\min(\dfrac{y}{1-\rho},L)}^{\min(\dfrac{y}{1-\rho},M)} \pi(x) f(y-(1-\rho)x)\,dx\nonumber\\&+f(y)\int_M^{+\infty}\pi(x)\,dx
\end{align*}

We can solve this last equation by fixed-point iteration algorithm by considering that $\pi(x)$ is the solution of $\pi(y) = g(\pi(y))$, where $g$ is a continuous function. Thus, for all $y$, we approximate $\pi$ at iteration $k$ by
\begin{align*}
w_k(y) = \int_0^{\min(L,y)} w_{k-1}(x) f(y-x)\,dx &+  \int_{\min(\dfrac{y}{1-\rho},L)}^{\min(\dfrac{y}{1-\rho},M)} w_{k-1} (x) f(y-(1-\rho)x)\,dx\\&+ f(y)\int_M^{+\infty}w_{k-1}(x)\,dx
\end{align*}

}
%%%%%%%%%%%%%%%%%%%%%%%%%%%%%%%%%%%%%%%%%

\frame{
\frametitle{Stationary Law}
\begin{figure}
\includegraphics[angle=0,width=0.45\textwidth]{{StationaryLaw}}
\includegraphics[angle=0,width=0.45\textwidth]{{StationnayLaw_L}}
\caption {The steady-state distribution in function of the inspection interval $\tau$ (left) and the preventive threshold $L$ (right):
$M=5$, $\rho=0.8$, $\alpha=1$, $\beta=1$, $b=1$.}
\end{figure}
}
%%%%%%%%%%%%%%%%%%%%%%%%%%%%%%%%%%%%%%%%%
%%%%%%%%%%%%%%%%%%%%%%%%%%%%%%%%%%%%%%%%%

\frame{
\frametitle{Sensitivity analysis}
\begin{figure}
\includegraphics[angle=0,width=0.45\textwidth]{{Cout_L_50}}
\includegraphics[angle=0,width=0.45\textwidth]{{Cout_L_20}}
\caption {Maintenance cost in function of the preventive threshold $L$, with $C_p=\frac{C_c}{2}$ (left) and $C_p=\frac{C_c}{5}$ (right) and $\tau=0.8$: $M=5$, $\rho=0.8$, $\alpha=1$, $\beta=1$, $b=1$.}
\end{figure}
}
%%%%%%%%%%%%%%%%%%%%%%%%%%%%%%%%%%%%%%%%%

\frame{
\frametitle{Sensitivity analysis}
\begin{figure}
\includegraphics[angle=0,width=0.45\textwidth]{{Cout_L_20_rho_half}}
\includegraphics[angle=0,width=0.45\textwidth]{{Cout_L_10_rho_dot8}}
\caption {Maintenance cost in function of the preventive threshold $L$, with $\tau=0.5$ (left) and $\tau=0.8$ (right): $M=5$, $\alpha=1$, $\beta=1$, $b=1$.}
\end{figure}
}
%%%%%%%%%%%%%%%%%%%%%%%%%%%%%%%%%%%%%%%%%
\frame{
\frametitle{Estimation}
Two cases are considered.
\begin{itemize}
	\item The deterioration level before and after maintenance is known and $\rho\sim \beta(r,s)$.\\
Estimation by the maximisation of the likelihood of the gamma distributed increments $\{y_{t_j^-}-y_{t_{j-1}^+}, j=0,1,\ldots\}$.\\
The maintenance efficiency is estimated by moment method on  $$\frac{y_{t_j^-}-y_{t_{j-1}^+}}{y_{t_j^-}},\;\; j=1,2,\ldots.$$
	\item The deterioration level in observed only before maintenance.
	Estimation by the maximisation of the likelihood of the gamma distributed increments $\{y_{t_j^-}-(1-\rho)y_{t_{j-1}^+}, j=0,1,\ldots\}$.\\
	The maintenance efficiency is estimated by moment method on  $$\frac{y_{t_j^-}-(1-\rho)y_{t_{j-1}^+}}{y_{t_j^-}},\;\; j=1,2,\ldots.$$
\end{itemize}


}

%%%%%%%%%%%%%%%%%%%%%%%%%%%%%%%%%%%%%%%%%
\frame{
\frametitle{}
\vspace{2cm}
\begin{center}
{\huge  Thank you for your attention}
\end{center}
}




%oooooooooooooooooooooooooooooooooooooooooooooooooooooooooooooooooooooo
\end{document}


